\documentclass[conference]{IEEEtran}
\usepackage{cite}
\usepackage{amsmath,amssymb}
\usepackage{graphicx,booktabs}
\usepackage{algorithm,algpseudocode}
\usepackage[hidelinks]{hyperref}
\begin{document}
\title{Combining Instance Filtering and Attribute Weighting in Decision Tree and Na\"ive Bayes Hybrids}
\author{
\IEEEauthorblockN{Ikramul Hasan Moral}
\IEEEauthorblockA{Department of Computer Science and Engineering\\
United International University\\
Dhaka, Bangladesh\\
imoral223489@bscse.uiu.ac.bd}
\and
\IEEEauthorblockN{Dewan Md. Farid}
\IEEEauthorblockA{Professor, Department of Computer Science and Engineering\\
United International University\\
Dhaka, Bangladesh\\
dewanfarid@cse.uiu.ac.bd}}


\maketitle
\begin{abstract}
Na\"ive Bayes (NB) can filter training instances before a decision tree is
fitted, and a decision tree can select and weight attributes for NB. We
join these two steps into Algorithm~3 and run them in both orders, with
either NB or C4.5 as the final classifier. All methods are tested on ten
benchmark datasets with ten repetitions of stratified 10-fold
cross-validation. Each cleaning step is fitted on the training fold only,
and every test row is scored. The two combined tree methods reach 81.21\%
and 80.89\% mean accuracy, against 86.33\% for C4.5 without cleaning. The
combined NB methods reach 75.63\% and 76.35\%, against 75.93\% for plain NB
and 77.35\% for attribute weighting alone. Paired tests over the datasets
show no clear accuracy difference between the two orders. When the filter
is fitted on the full dataset before cross-validation, the mean tree
accuracy of Algorithm~1 rises by 13.50 percentage points, but the filter
also removes cases that would later be tested. The removal records show
that the filter can drop a much larger share of one class than of another.
On these datasets we found no evidence that combining the steps improves
on the base classifiers, and the scores cannot be read correctly without
knowing where cleaning took place and which cases it removed.
\end{abstract}
\begin{IEEEkeywords}
Na\"ive Bayes, decision tree, instance filtering, attribute weighting,
hybrid classifier, cross-validation
\end{IEEEkeywords}

\section{Introduction}
Farid et al.~\cite{farid2014hybrid} introduced two hybrids that connect
Na\"ive Bayes (NB) and decision trees. Algorithm~1 trains NB on the training
set, removes the instances it misclassifies, and grows a tree on the
remaining rows. Algorithm~2 grows a tree first, keeps the attributes it
tests, and gives more weight to attributes tested closer to the root when
fitting NB. In each method, one classifier prepares the data for the other.

Running both steps is the obvious extension, but combining them changes
the data each learner sees. If instances are removed first, the
tree picks attributes from a smaller sample. If attributes are picked
first, a weighted NB judge decides which instances stay. We call the
combined procedure Algorithm~3 and test both orders. We report a final
tree and a final NB separately, so that a gain from switching classifiers
is not mistaken for a gain from cleaning.

Where the cleaning happens also changes what a score means. A filter fitted
before cross-validation can see the labels of rows that will later be test
rows, and it can remove those rows before they are ever scored. A filter
fitted inside each training fold leaves every test row in place. The first
approach therefore changes both the information the cleaner uses and the
set of rows on which accuracy is measured.

We compare Algorithm~3 with each single step and with the uncleaned
baselines, using the same fold-specific protocol for all methods. We report
accuracy, macro-F1, and the share of each class removed by the filter, and
we compare cleaning inside the training fold with cleaning before the
split. On this benchmark the combined methods show no clear advantage. The
contribution of this paper is a controlled comparison, together with an
account of how the evaluation protocol and the removed rows affect the
reported scores.

\section{Related work}
Hall~\cite{hall2007decision} used decision trees to derive attribute weights
for NB as a way to reduce the effect of its conditional-independence
assumption. Farid et al.~\cite{farid2014hybrid} proposed the two steps
studied here and evaluated each one on its own. Algorithm~3 chains those
steps and keeps the same tree learner and NB likelihood model.

Instance filtering needs more care. Brodley and
Friedl~\cite{brodley1999identifying} used several classifiers to identify
mislabeled training data. When a single NB judge disagrees with a label,
that alone does not show the label is wrong. In this paper, a
``filtered'' row is one the judge rejected. We make no claim that
the removed rows contain real label errors.

Cawley and Talbot~\cite{cawley2010overfitting} showed that evaluation
choices can introduce selection bias large enough to affect comparisons
between learning algorithms. We ask the same kind of question about where
supervised cleaning is placed, for one family of hybrids. Following
Dem\v{s}ar~\cite{demsar2006statistical}, we make paired comparisons at the
dataset level. The repeated folds estimate each dataset's score and
are not treated as independent samples.

\section{Algorithms}
\subsection{Algorithm 1: instance filtering}
Given training data $D$, the NB judge is fitted on $D$ and predicts the
labels of those same rows. A row is kept only if the prediction matches its
given label. Algorithm~1 then fits C4.5 on the kept rows. The reference arm
that ends in NB uses the same filter but fits NB again as the final
classifier. If deletion would leave fewer than two classes, our
implementation keeps the original training rows.

This is a disagreement rule on the training set, so its predictions are not
out-of-fold estimates of label quality. A judge can reject a correctly
labeled case when the classes overlap or when its own model assumptions do
not fit the data.

\subsection{Algorithm 2: attribute selection and weighting}
A C4.5 tree is fitted on the training data. Let $d_j$ be the smallest depth
at which it tests attribute $j$, with the root at depth one. The weight is
\begin{equation}
 w_j = \begin{cases}1/\sqrt{d_j}, & j\text{ is tested},\\
                    0, & \text{otherwise}.\end{cases}
\end{equation}
Attributes the tree does not test are dropped. The weighted NB class score is
\begin{equation}
 s_c(x)=\log P(c)+\sum_{j\in S}w_j\log P(x_j\mid c),
\end{equation}
where $S$ is the set of kept attributes. NB predicts the class with the
largest score. The reference arm that ends in a tree keeps only the
selected columns, and its final C4.5 does not use the weights.

\subsection{Algorithm 3: the combined hybrid}
The instances-first path, NA, first runs the plain NB filter of
Algorithm~1. It then fits a tree on the kept rows to obtain the selected
attributes and their weights. The attributes-first path, AN, takes the
attributes and weights from the full training fold and then uses weighted
NB to filter the rows. In AN, the judge and the final NB use the same
weights, and the weights are not estimated again after deletion. Each path
ends in either weighted NB or C4.5.

\setcounter{algorithm}{2}
\begin{algorithm}[t]
\caption{Combined cleaning and classification}
\label{alg:three}
\begin{algorithmic}[1]
\Require Training data $D$, order $o\in\{\mathrm{NA},\mathrm{AN}\}$,
 final learner $L\in\{\mathrm{NB},\mathrm{DT}\}$
\If{$o=\mathrm{NA}$}
 \State $D\gets\Call{Filter}{D,\text{plain NB}}$
 \State $(S,W)\gets\Call{Select}{D,\text{C4.5}}$
\Else
 \State $(S,W)\gets\Call{Select}{D,\text{C4.5}}$
 \State $D\gets D[:,S]$
 \State $D\gets\Call{Filter}{D,\text{NB with }W}$
\EndIf
\If{$o=\mathrm{NA}$}
 \State $D\gets D[:,S]$
\EndIf
\State Fit $M$ on $D$ using $L$ (use $W$ if $L=\mathrm{NB}$)
\State \Return $M,S$
\end{algorithmic}
\end{algorithm}

Algorithm~\ref{alg:three} uses the same filtering safeguard as
Algorithm~1. If the selection tree tests no attribute, the combined path
keeps all current columns with unit weights. This never happened in the
stored runs. At prediction time, the test fold keeps only the columns in
$S$ but all of its rows, so every test case is classified.

\section{Experimental method}
\subsection{Datasets and implementation}
We use the ten datasets named in the 2014 study. Table~\ref{tab:data}
lists the sizes of the copies we used. Image segmentation joins the two
UCI files, and NSL-KDD uses its 20\% training file
\cite{tavallaee2009detailed}. Our Glass and NSL-KDD files each have one
class fewer than the original study lists. The data come from
the same benchmark family, but this is not an exact reproduction of the
original data or splits.

\begin{table}[t]
\caption{Dataset sizes in the files used for evaluation}
\label{tab:data}
\centering
\begin{tabular}{lrrr}
\toprule
Dataset & Rows & Attributes & Classes \\
\midrule
Breast cancer & 286 & 9 & 2 \\
Contact lenses & 24 & 4 & 3 \\
Diabetes & 768 & 8 & 2 \\
Glass & 214 & 9 & 6 \\
Iris & 150 & 4 & 3 \\
Soybean & 683 & 35 & 19 \\
Vote & 435 & 16 & 2 \\
Image segmentation & 2310 & 19 & 7 \\
Tic-tac-toe & 958 & 9 & 2 \\
NSL-KDD & 25192 & 41 & 22 \\
\bottomrule
\end{tabular}
\end{table}

C4.5 is Weka J48 version 3.8.6 with pruning enabled and the default
settings \texttt{-C 0.25 -M 2}~\cite{witten2016weka}. Our NB uses add-one
counts for nominal attributes and Gaussian likelihoods for numeric
attributes. Nominal attributes stay as single columns, so a tree selects an
original attribute rather than separate one-hot indicators. The category
list of each nominal attribute comes from the dataset file and is fixed
across folds. A missing nominal value is treated as its own category.
Counts, numeric distributions, and class priors are estimated from the
training rows alone.

\subsection{Protocols and comparison arms}
Protocol B is our main evaluation. We run stratified 10-fold
cross-validation with seeds 0--9, which gives 100 scored folds per dataset
and method. All methods use the same splits. The judge, the selection tree,
and the final classifier are all fitted on the training rows of each split.
Only the choice of columns is carried over to the test fold, and test rows
are never filtered. This is repeated cross-validation with preprocessing
inside each fold. It is not a nested search over hyperparameters.

Protocol A fits the cleaning steps once on the whole dataset and then runs
cross-validation for the final classifier on the cleaned data. The cleaner
therefore sees the labels of future test rows. For methods that filter
instances, it also removes future test rows, so the folds are built from a
different sample. We report these score changes as descriptive results
only. Protocol A folds are not paired with Protocol B folds, and the
difference between the two cannot separate label leakage from the change in
which rows are evaluated.

The comparison includes plain NB and C4.5, the original Algorithm~1--DT
and Algorithm~2--NB methods, and the four versions of Algorithm~3. Two
reference arms help show where changes come from. Algorithm~1--NB changes
only the final learner after filtering, and Algorithm~2--DT applies
attribute selection before another tree. In the tables, N stands for
filtering, A for attribute selection and weighting, and NA and AN for the
two combined orders.

\subsection{Measures and paired analysis}
For each method and dataset, accuracy and macro-F1 are averaged over the
100 scored folds, with equal weight for each fold. Within a fold, macro-F1
averages the F1 scores of the classes that appear in the true labels or the
predictions, and a class with no correct prediction counts as zero. The
summary figures then give each of the ten datasets equal weight, so they
are not pooled scores over all instances.

We compare eight pairs with two-sided Wilcoxon signed-rank tests on the
ten dataset-level mean accuracies. The pairs are the two original hybrids
against their baselines, each of the four combined methods against its
baseline, and the two order comparisons. We compute exact $p$ values by
enumerating all sign permutations, which suits the small sample and handles
zero differences. Holm adjustment controls the family of eight tests at
0.05. We chose these tests after the experiments had been run and applied
them to the saved results, so they should be read as exploratory. Mean
differences are given in percentage points (pp). Differences between the
combined methods and the single-step methods are reported as means only,
without tests. A nonsignificant result does not show that two methods are
equivalent.

For each filtering step we also record the share of each class that is
removed, using that class's number of training rows as the denominator.
This shows how deletion changes the class balance without assuming that
the removed labels were wrong.

\section{Results and discussion}
\subsection{Combined methods have lower mean tree accuracy}
Table~\ref{tab:results} gives the averages over the benchmark, and
Table~\ref{tab:results-datasets} gives the results for each dataset.
Uncleaned C4.5 reaches 86.33\% accuracy and 81.92\% macro-F1. The two
combined tree methods reach 81.21\% and 80.89\% accuracy, which is 5.12 and
5.44 pp lower. Their macro-F1 also drops, to 75.76\% and 73.93\%. Compared
with the 81.03\% of Algorithm~1, adding attribute selection after filtering
gains only 0.18 pp.

\begin{table}[t]
\caption{Protocol B results (\%), averaged equally across ten datasets}
\label{tab:results}
\centering
% generated by scripts/make_results_table.py; do not edit
% 10 datasets, seeds 0..9,
% Weka 3.8.6 pruned, defaults -C 0.25 -M 2, FaithfulNB (add-one counts for nominal, Gaussian for numeric)
\begin{tabular}{lrr}
\toprule
Method & Accuracy & Macro-F1 \\
\midrule
NB & 75.93 & 69.80 \\
C4.5 & 86.33 & 81.92 \\
Algorithm 1 & 81.03 & 75.70 \\
Algorithm 2 & 77.35 & 71.73 \\
Instances first, NB & 75.63 & 69.60 \\
Instances first, DT & 81.21 & 75.76 \\
Attributes first, NB & 76.35 & 69.91 \\
Attributes first, DT & 80.89 & 73.93 \\
\bottomrule
\end{tabular}

\end{table}

\begin{table*}[t]
\caption{Protocol B accuracy by dataset (\%). Base is the uncleaned
classifier; N is instance filtering; A is attribute selection and weighting;
NA and AN are the two Algorithm~3 orders. N--NB and A--DT are reference arms.}
\label{tab:results-datasets}
\centering
\footnotesize
\setlength{\tabcolsep}{3pt}
\begin{tabular}{lrrrrrrrrrr}
\toprule
 & \multicolumn{5}{c}{Final NB} & \multicolumn{5}{c}{Final C4.5} \\
\cmidrule(lr){2-6}\cmidrule(lr){7-11}
Dataset & Base & N & A & NA & AN & Base & N & A & NA & AN \\
\midrule
Breast cancer & 72.12 & 72.68 & 72.43 & 71.28 & 72.53 & 74.84 & 71.77 & 73.91 & 71.77 & 72.08 \\
Contact lenses & 73.00 & 73.00 & 80.67 & 83.00 & 79.33 & 84.67 & 84.67 & 84.67 & 84.67 & 80.67 \\
Diabetes & 75.52 & 74.41 & 76.73 & 75.58 & 75.67 & 74.12 & 75.12 & 74.18 & 75.23 & 75.57 \\
Glass & 45.95 & 37.00 & 47.84 & 30.93 & 37.28 & 66.58 & 48.83 & 66.58 & 48.97 & 49.15 \\
Segmentation & 79.76 & 74.92 & 82.30 & 74.53 & 77.32 & 96.81 & 89.24 & 96.83 & 89.26 & 90.88 \\
Iris & 95.53 & 95.33 & 95.40 & 95.33 & 95.87 & 94.27 & 94.93 & 94.27 & 94.93 & 93.80 \\
NSL-KDD & 67.42 & 75.79 & 62.80 & 74.53 & 75.23 & 99.47 & 87.43 & 99.47 & 89.68 & 87.07 \\
Soybean & 90.09 & 89.00 & 89.40 & 85.67 & 86.66 & 92.71 & 88.71 & 92.26 & 87.89 & 90.98 \\
Tic-tac-toe & 69.83 & 71.89 & 70.68 & 70.03 & 68.07 & 84.74 & 74.43 & 84.74 & 74.55 & 73.17 \\
Vote & 90.06 & 90.13 & 95.24 & 95.45 & 95.49 & 95.08 & 95.13 & 95.03 & 95.13 & 95.54 \\
\bottomrule
\end{tabular}

\end{table*}

For NB, attribute weighting alone has the highest mean accuracy, 77.35\%
against 75.93\% for plain NB. The combined methods reach 75.63\% and
76.35\%, or $-0.30$ and $+0.42$ pp relative to plain NB, and neither one
matches Algorithm~2. The reference arms point to filtering as the main
source of the loss. Filtering followed by NB reaches 75.42\%, while
selecting attributes before C4.5 leaves its mean almost unchanged at
86.19\%.

The averages hide large differences between datasets. On Glass,
instances-first NB falls from 45.95\% for plain NB to 30.93\%, and
instances-first C4.5 falls from 66.58\% to 48.97\%. On Contact lenses, the
same NB combination rises from 73.00\% to 83.00\%, but that dataset has
only 24 rows. A cleaner should be checked dataset by dataset before it is
called helpful.

Table~\ref{tab:tests} reports the paired accuracy tests. AN--DT is below
the C4.5 baseline with raw $p=0.014$, but its Holm-adjusted $p=0.109$ is
above the threshold. None of the eight comparisons is significant after
adjustment. The mean losses are still worth noting, but this benchmark does
not show a general advantage or disadvantage across classification
problems.

\begin{table}[t]
\caption{Dataset-level paired accuracy tests. Differences are first method
minus second; Holm adjustment covers all eight comparisons.}
\label{tab:tests}
\centering
\footnotesize
\begin{tabular}{lrrr}
\toprule
Comparison & $\Delta$ (pp) & $p$ & Holm $p$ \\
\midrule
N--DT vs. DT & -5.30 & 0.055 & 0.383 \\
A--NB vs. NB & +1.42 & 0.131 & 0.654 \\
NA--DT vs. DT & -5.12 & 0.055 & 0.383 \\
AN--DT vs. DT & -5.44 & 0.014 & 0.109 \\
NA--NB vs. NB & -0.30 & 1.000 & 1.000 \\
AN--NB vs. NB & +0.42 & 0.846 & 1.000 \\
AN--DT vs. NA--DT & -0.32 & 0.922 & 1.000 \\
AN--NB vs. NA--NB & +0.71 & 0.275 & 1.000 \\
\bottomrule
\end{tabular}

\end{table}

Switching from NA to AN raises mean NB accuracy by 0.71 pp and lowers mean
tree accuracy by 0.32 pp, with raw $p$ values of 0.275 and 0.922. We see no
clear order effect under these settings. This result is limited to a single
pass through the two steps, with fixed learner settings and with the AN
judge using the selection weights.

\subsection{Cleaning before the split changes the score}
Table~\ref{tab:results-protocol-a} compares the two protocols. Algorithm~1
reaches 94.52\% when the filter is fitted on the full data, against
81.03\% when the judge is fitted inside each training fold. That is an
apparent gain of 13.50 pp. Algorithm~2 changes much less, from 77.35\% to
78.07\%. The combined tree methods also jump under Protocol A, to 94.72\%
and 96.32\%.

\begin{table}[t]
\caption{Mean accuracy (\%) under the two protocols. Filtering before CV
changes which rows are evaluated, so these scores should be read as a
diagnostic comparison.}
\label{tab:results-protocol-a}
\centering
\footnotesize
% generated by scripts/make_results_table.py; do not edit
% 10 datasets, seeds 0..9,
% Weka 3.8.6 pruned, defaults -C 0.25 -M 2, FaithfulNB (add-one counts for nominal, Gaussian for numeric)
\begin{tabular}{lrrr}
\toprule
Method & Fold-specific & Before CV & Difference \\
\midrule
Algorithm 1 & 81.03 & 94.52 & +13.50 \\
Algorithm 2 & 77.35 & 78.07 & +0.72 \\
Alg.\ 1, NB final & 75.42 & 91.93 & +16.51 \\
Alg.\ 2, DT final & 86.19 & 86.52 & +0.33 \\
Instances first, NB & 75.63 & 90.11 & +14.48 \\
Instances first, DT & 81.21 & 94.72 & +13.51 \\
Attributes first, NB & 76.35 & 91.08 & +14.74 \\
Attributes first, DT & 80.89 & 96.32 & +15.43 \\
\bottomrule
\end{tabular}

\end{table}

The two protocols measure different things. Protocol B asks the final
classifier to predict every held-out case. Protocol A first lets the NB
judge decide which cases reach the test folds. The large increase fits
with the judge removing cases it finds hard, even if the final classifier
does no better on the full data. Attribute-only cleaning keeps all rows,
but selecting attributes on the full data still uses information from
future test folds. Its small score change here does not mean that
attributes can safely be selected before cross-validation.

\subsection{Deletion changes class balance}
The breast-cancer records show how uneven the removal can be. Across
Protocol B folds, the plain NB judge removes on average 50.63\% of the
recurrence-event rows but only 14.51\% of the no-recurrence-event rows.
After attribute selection, the weighted judge removes 58.97\% and 11.72\%.
These are mean removal rates within each class, not shares of the whole
dataset.

In both paths the minority class loses a much larger share of its training
rows. The records cannot tell us whether those labels were wrong or whether
NB simply failed on hard but valid cases. They do show that a cleaner
should be judged by the rows it removes as well as by its scores.

\section{Limitations}
The benchmark uses the dataset family of the original study, with the file
differences noted above. It says nothing about performance on modern
datasets or on data with a known level of label noise. In particular, our
NSL-KDD C4.5 baseline is 99.47\%, while the 2014 study reports 71.11\%. We
have not been able to explain this gap. The comparison here uses shared
data and splits, but it should not be described as an exact replication of
the earlier scores.

Rare classes also limit the evaluation. Contact lenses has classes with
four and five rows, and NSL-KDD has several classes with fewer than ten
examples, four of which have a single row. As a result, some stratified folds
miss a class in training or testing. Repeating the splits does not create
new examples, and the per-fold macro-F1 can average over different class
sets. These results do not support reliable claims about rare classes.

In Protocol B, model fitting and supervised cleaning use only training
rows, but the category lists of nominal attributes come from the complete
files. The estimates therefore assume these categories are known in
advance, and they do not test unseen nominal values. We did not use
calibration or out-of-fold judging to check the rejection decisions of the
filter. With ten datasets, the tests can miss small effects. Finally, the
Protocol A results cannot tell us which evaluation procedure the 2014
study used.

\section{Conclusion}
Algorithm~3 combines NB instance filtering and tree-based attribute
weighting in either order. On ten datasets with fold-specific evaluation,
both combined tree methods have lower mean accuracy and macro-F1 than
uncleaned C4.5, and both combined NB methods fall short of attribute
weighting alone. The paired tests show no clear order effect. The
before-CV scores and the removal rates per class lead to a practical
recommendation. Supervised cleaning should be fitted inside the training
fold so that every held-out row is still scored, and studies that use such
a filter should report how many training cases each class loses.

\bibliographystyle{IEEEtran}
\bibliography{references}
\end{document}
